\appendix
\chapter{Pregled zaglavlja \texttt{<algorithm>} (C++11)}
Ovaj dodatak daje pregled algoritama,
odnosno generičkih funkcija koje uglavnom koriste iteratore i funktore,
a koje su definirane u C++11 zaglavlju \texttt{<algorithm>}.
Za svaku funkciju navedeno je njeno ime, kratak opis
te jedan jednostavan primjer njene upotrebe.

\begin{itemize}
\item \textbf{\texttt{adjacent\_find}}\newline
Pronalazi prvi par susjednih jednakih elemenata.
\begin{minted}{C++}
auto it = std::adjacent_find(v.begin(), v.end());
\end{minted}

\item \textbf{\texttt{all\_of}}\newline
Provjerava zadovoljavaju li svi elementi uvjet.
\begin{minted}{C++}
bool b = std::all_of(v.begin(), v.end(), [](int x){ return x > 0; });
\end{minted}

\item \textbf{\texttt{any\_of}}\newline
Provjerava postoji li barem jedan element koji zadovoljava uvjet.
\begin{minted}{C++}
bool b = std::any_of(v.begin(), v.end(), [](int x){ return x < 0; });
\end{minted}

\item \textbf{\texttt{binary\_search}}\newline
Provjerava postoji li element u sortiranom rasponu.
\begin{minted}{C++}
bool b = std::binary_search(v.begin(), v.end(), 10);
\end{minted}

\item \textbf{\texttt{copy}}\newline
Kopira elemente u drugi raspon.
\begin{minted}{C++}
std::copy(v.begin(), v.end(), u.begin());
\end{minted}

\item \textbf{\texttt{copy\_backward}}\newline
Kopira elemente unatrag.
\begin{minted}{C++}
std::copy_backward(v.begin(), v.end(), u.end());
\end{minted}

\item \textbf{\texttt{copy\_if}}\newline
Kopira samo elemente koji zadovoljavaju uvjet.
\begin{minted}{C++}
std::copy_if(v.begin(), v.end(), std::back_inserter(u), [](int x){ return x > 0; });
\end{minted}

\item \textbf{\texttt{copy\_n}}\newline
Kopira prvih n elemenata.
\begin{minted}{C++}
std::copy_n(v.begin(), 5, u.begin());
\end{minted}

\item \textbf{\texttt{count}}\newline
Broji pojavljivanja vrijednosti.
\begin{minted}{C++}
int n = std::count(v.begin(), v.end(), 5);
\end{minted}

\item \textbf{\texttt{count\_if}}\newline
Broji elemente koji zadovoljavaju uvjet.
\begin{minted}{C++}
int n = std::count_if(v.begin(), v.end(), [](int x){ return x % 2 == 0; });
\end{minted}

\item \textbf{\texttt{equal}}\newline
Provjerava jednakost dvaju raspona.
\begin{minted}{C++}
bool b = std::equal(v.begin(), v.end(), u.begin());
\end{minted}

\item \textbf{\texttt{equal\_range}}\newline
Daje raspon jednakih elemenata u sortiranom nizu.
\begin{minted}{C++}
auto p = std::equal_range(v.begin(), v.end(), 5);
\end{minted}

\item \textbf{\texttt{fill}}\newline
Popunjava raspon vrijednošću.
\begin{minted}{C++}
std::fill(v.begin(), v.end(), 0);
\end{minted}

\item \textbf{\texttt{fill\_n}}\newline
Popunjava prvih n elemenata.
\begin{minted}{C++}
std::fill_n(v.begin(), 10, 0);
\end{minted}

\item \textbf{\texttt{find}}\newline
Pronalazi prvi element zadane vrijednosti.
\begin{minted}{C++}
auto it = std::find(v.begin(), v.end(), 5);
\end{minted}

\item \textbf{\texttt{find\_end}}\newline
Pronalazi zadnje pojavljivanje podraspona.
\begin{minted}{C++}
auto it = std::find_end(v.begin(), v.end(), p.begin(), p.end());
\end{minted}

\item \textbf{\texttt{find\_first\_of}}\newline
Pronalazi prvi element koji pripada drugom rasponu.
\begin{minted}{C++}
auto it = std::find_first_of(v.begin(), v.end(), p.begin(), p.end());
\end{minted}

\item \textbf{\texttt{find\_if}}\newline
Pronalazi prvi element koji zadovoljava uvjet.
\begin{minted}{C++}
auto it = std::find_if(v.begin(), v.end(), [](int x){ return x > 10; });
\end{minted}

\item \textbf{\texttt{find\_if\_not}}\newline
Pronalazi prvi element koji ne zadovoljava uvjet.
\begin{minted}{C++}
auto it = std::find_if_not(v.begin(), v.end(), [](int x){ return x > 0; });
\end{minted}

\item \textbf{\texttt{for\_each}}\newline
Primjenjuje funkciju na sve elemente.
\begin{minted}{C++}
std::for_each(v.begin(), v.end(), [](int x){ std::cout << x; });
\end{minted}

\item \textbf{\texttt{generate}}\newline
Generira vrijednosti pomoću funkcije.
\begin{minted}{C++}
std::generate(v.begin(), v.end(), gen);
\end{minted}

\item \textbf{\texttt{generate\_n}}\newline
Generira prvih n vrijednosti.
\begin{minted}{C++}
std::generate_n(v.begin(), 5, gen);
\end{minted}

\item \textbf{\texttt{includes}}\newline
Provjerava sadrži li raspon drugi sortiran raspon.
\begin{minted}{C++}
bool b = std::includes(a.begin(), a.end(), b.begin(), b.end());
\end{minted}

\item \textbf{\texttt{inplace\_merge}}\newline
Spaja dva sortirana dijela.
\begin{minted}{C++}
std::inplace_merge(v.begin(), mid, v.end());
\end{minted}

\item \textbf{\texttt{is\_heap}}\newline
Provjerava je li raspon gomila (heap).
\begin{minted}{C++}
bool b = std::is_heap(v.begin(), v.end());
\end{minted}

\item \textbf{\texttt{is\_heap\_until}}\newline
Pronalazi kraj heap strukture.
\begin{minted}{C++}
auto it = std::is_heap_until(v.begin(), v.end());
\end{minted}

\item \textbf{\texttt{is\_partitioned}}\newline
Provjerava je li raspon particioniran.
\begin{minted}{C++}
bool b = std::is_partitioned(v.begin(), v.end(), pred);
\end{minted}

\item \textbf{\texttt{is\_permutation}}\newline
Provjerava je li permutacija drugog raspona.
\begin{minted}{C++}
bool b = std::is_permutation(v.begin(), v.end(), u.begin());
\end{minted}

\item \textbf{\texttt{is\_sorted}}\newline
Provjerava je li sortiran.
\begin{minted}{C++}
bool b = std::is_sorted(v.begin(), v.end());
\end{minted}

\item \textbf{\texttt{is\_sorted\_until}}\newline
Pronalazi kraj sortiranog dijela.
\begin{minted}{C++}
auto it = std::is_sorted_until(v.begin(), v.end());
\end{minted}

\item \textbf{\texttt{iter\_swap}}\newline
Zamjenjuje elemente na iteratorima.
\begin{minted}{C++}
std::iter_swap(it1, it2);
\end{minted}

\item \textbf{\texttt{lexicographical\_compare}}\newline
Leksičko uspoređivanje raspona.
\begin{minted}{C++}
bool b = std::lexicographical_compare(a.begin(), a.end(), b.begin(), b.end());
\end{minted}

\item \textbf{\texttt{lower\_bound}}\newline
Pronalazi prvu poziciju elementa.
\begin{minted}{C++}
auto it = std::lower_bound(v.begin(), v.end(), 10);
\end{minted}

\item \textbf{\texttt{make\_heap}}\newline
Stvara heap strukturu.
\begin{minted}{C++}
std::make_heap(v.begin(), v.end());
\end{minted}

\item \textbf{\texttt{max}}\newline
Vraća veći od dva elementa.
\begin{minted}{C++}
auto m = std::max(a, b);
\end{minted}

\item \textbf{\texttt{max\_element}}\newline
Pronalazi najveći element.
\begin{minted}{C++}
auto it = std::max_element(v.begin(), v.end());
\end{minted}

\item \textbf{\texttt{merge}}\newline
Spaja dva sortirana raspona.
\begin{minted}{C++}
std::merge(a.begin(), a.end(), b.begin(), b.end(), out.begin());
\end{minted}

\item \textbf{\texttt{min}}\newline
Vraća manji od dva elementa.
\begin{minted}{C++}
auto m = std::min(a, b);
\end{minted}

\item \textbf{\texttt{min\_element}}\newline
Pronalazi najmanji element.
\begin{minted}{C++}
auto it = std::min_element(v.begin(), v.end());
\end{minted}

\item \textbf{\texttt{minmax}}\newline
Vraća i min i max.
\begin{minted}{C++}
auto p = std::minmax(a, b);
\end{minted}

\item \textbf{\texttt{minmax\_element}}\newline
Pronalazi min i max element.
\begin{minted}{C++}
auto p = std::minmax_element(v.begin(), v.end());
\end{minted}

\item \textbf{\texttt{mismatch}}\newline
Pronalazi prvo neslaganje dvaju raspona.
\begin{minted}{C++}
auto p = std::mismatch(v.begin(), v.end(), u.begin());
\end{minted}

\item \textbf{\texttt{move}}\newline
Premješta elemente.
\begin{minted}{C++}
std::move(v.begin(), v.end(), u.begin());
\end{minted}

\item \textbf{\texttt{move\_backward}}\newline
Premješta elemente unatrag.
\begin{minted}{C++}
std::move_backward(v.begin(), v.end(), u.end());
\end{minted}

\item \textbf{\texttt{next\_permutation}}\newline
Generira sljedeću permutaciju.
\begin{minted}{C++}
std::next_permutation(v.begin(), v.end());
\end{minted}

\item \textbf{\texttt{none\_of}}\newline
Provjerava da nijedan element ne zadovoljava uvjet.
\begin{minted}{C++}
bool b = std::none_of(v.begin(), v.end(), pred);
\end{minted}

\item \textbf{\texttt{nth\_element}}\newline
Postavlja n-ti element na ispravno mjesto.
\begin{minted}{C++}
std::nth_element(v.begin(), v.begin()+k, v.end());
\end{minted}

\item \textbf{\texttt{partial\_sort}}\newline
Djelomično sortira raspon.
\begin{minted}{C++}
std::partial_sort(v.begin(), v.begin()+k, v.end());
\end{minted}

\item \textbf{\texttt{partial\_sort\_copy}}\newline
Kopira i sortira dio raspona.
\begin{minted}{C++}
std::partial_sort_copy(v.begin(), v.end(), u.begin(), u.end());
\end{minted}

\item \textbf{\texttt{partition}}\newline
Dijeli raspon prema uvjetu.
\begin{minted}{C++}
std::partition(v.begin(), v.end(), pred);
\end{minted}

\item \textbf{\texttt{partition\_copy}}\newline
Kopira u dvije skupine.
\begin{minted}{C++}
std::partition_copy(v.begin(), v.end(), a.begin(), b.begin(), pred);
\end{minted}

\item \textbf{\texttt{partition\_point}}\newline
Pronalazi granicu particije.
\begin{minted}{C++}
auto it = std::partition_point(v.begin(), v.end(), pred);
\end{minted}

\item \textbf{\texttt{pop\_heap}}\newline
Uklanja element iz heap-a.
\begin{minted}{C++}
std::pop_heap(v.begin(), v.end());
\end{minted}

\item \textbf{\texttt{prev\_permutation}}\newline
Generira prethodnu permutaciju.
\begin{minted}{C++}
std::prev_permutation(v.begin(), v.end());
\end{minted}

\item \textbf{\texttt{random\_shuffle}}\newline
Nasumično miješa elemente.
\begin{minted}{C++}
std::random_shuffle(v.begin(), v.end());
\end{minted}

\item \textbf{\texttt{remove}}\newline
Uklanja elemente logički.
\begin{minted}{C++}
v.erase(std::remove(v.begin(), v.end(), 0), v.end());
\end{minted}

\item \textbf{\texttt{remove\_copy}}\newline
Kopira bez određenih elemenata.
\begin{minted}{C++}
std::remove_copy(v.begin(), v.end(), u.begin(), 0);
\end{minted}

\item \textbf{\texttt{remove\_copy\_if}}\newline
Kopira uz uvjet.
\begin{minted}{C++}
std::remove_copy_if(v.begin(), v.end(), u.begin(), pred);
\end{minted}

\item \textbf{\texttt{remove\_if}}\newline
Uklanja elemente po uvjetu.
\begin{minted}{C++}
v.erase(std::remove_if(v.begin(), v.end(), pred), v.end());
\end{minted}

\item \textbf{\texttt{replace}}\newline
Zamjenjuje vrijednosti.
\begin{minted}{C++}
std::replace(v.begin(), v.end(), 0, 1);
\end{minted}

\item \textbf{\texttt{replace\_copy}}\newline
Kopira uz zamjenu.
\begin{minted}{C++}
std::replace_copy(v.begin(), v.end(), u.begin(), 0, 1);
\end{minted}

\item \textbf{\texttt{replace\_copy\_if}}\newline
Kopira uz uvjetnu zamjenu.
\begin{minted}{C++}
std::replace_copy_if(v.begin(), v.end(), u.begin(), pred, 1);
\end{minted}

\item \textbf{\texttt{replace\_if}}\newline
Zamjenjuje po uvjetu.
\begin{minted}{C++}
std::replace_if(v.begin(), v.end(), pred, 1);
\end{minted}

\item \textbf{\texttt{reverse}}\newline
Obrće redoslijed.
\begin{minted}{C++}
std::reverse(v.begin(), v.end());
\end{minted}

\item \textbf{\texttt{reverse\_copy}}\newline
Kopira obrnuto.
\begin{minted}{C++}
std::reverse_copy(v.begin(), v.end(), u.begin());
\end{minted}

\item \textbf{\texttt{rotate}}\newline
Rotira elemente.
\begin{minted}{C++}
std::rotate(v.begin(), v.begin()+k, v.end());
\end{minted}

\item \textbf{\texttt{rotate\_copy}}\newline
Kopira rotirano.
\begin{minted}{C++}
std::rotate_copy(v.begin(), v.begin()+k, v.end(), u.begin());
\end{minted}

\item \textbf{\texttt{search}}\newline
Pretražuje podraspon.
\begin{minted}{C++}
auto it = std::search(v.begin(), v.end(), p.begin(), p.end());
\end{minted}

\item \textbf{\texttt{search\_n}}\newline
Traži n uzastopnih elemenata.
\begin{minted}{C++}
auto it = std::search_n(v.begin(), v.end(), 3, 5);
\end{minted}

\item \textbf{\texttt{set\_difference}}\newline
Razlika skupova.
\begin{minted}{C++}
std::set_difference(a.begin(), a.end(), b.begin(), b.end(), out.begin());
\end{minted}

\item \textbf{\texttt{set\_intersection}}\newline
Presjek skupova.
\begin{minted}{C++}
std::set_intersection(a.begin(), a.end(), b.begin(), b.end(), out.begin());
\end{minted}

\item \textbf{\texttt{set\_symmetric\_difference}}\newline
Simetrična razlika.
\begin{minted}{C++}
std::set_symmetric_difference(a.begin(), a.end(), b.begin(), b.end(), out.begin());
\end{minted}

\item \textbf{\texttt{set\_union}}\newline
Unija skupova.
\begin{minted}{C++}
std::set_union(a.begin(), a.end(), b.begin(), b.end(), out.begin());
\end{minted}

\item \textbf{\texttt{shuffle}}\newline
Nasumično miješanje.
\begin{minted}{C++}
std::shuffle(v.begin(), v.end(), gen);
\end{minted}

\item \textbf{\texttt{sort}}\newline
Sortira elemente.
\begin{minted}{C++}
std::sort(v.begin(), v.end());
\end{minted}

\item \textbf{\texttt{stable\_partition}}\newline
Stabilna particija.
\begin{minted}{C++}
std::stable_partition(v.begin(), v.end(), pred);
\end{minted}

\item \textbf{\texttt{stable\_sort}}\newline
Stabilno sortiranje.
\begin{minted}{C++}
std::stable_sort(v.begin(), v.end());
\end{minted}

\item \textbf{\texttt{swap}}\newline
Zamjena vrijednosti.
\begin{minted}{C++}
std::swap(a, b);
\end{minted}

\item \textbf{\texttt{swap\_ranges}}\newline
Zamjena raspona.
\begin{minted}{C++}
std::swap_ranges(a.begin(), a.end(), b.begin());
\end{minted}

\item \textbf{\texttt{transform}}\newline
Transformira elemente.
\begin{minted}{C++}
std::transform(v.begin(), v.end(), v.begin(), op);
\end{minted}

\item \textbf{\texttt{unique}}\newline
Uklanja uzastopne duplikate.
\begin{minted}{C++}
v.erase(std::unique(v.begin(), v.end()), v.end());
\end{minted}

\item \textbf{\texttt{unique\_copy}}\newline
Kopira bez duplikata.
\begin{minted}{C++}
std::unique_copy(v.begin(), v.end(), u.begin());
\end{minted}

\item \textbf{\texttt{upper\_bound}}\newline
Pronalazi gornju granicu.
\begin{minted}{C++}
auto it = std::upper_bound(v.begin(), v.end(), 10);
\end{minted}
\end{itemize}